%% kapa: exam=1 round=23 year=2012 period=1 form=A subject=civil_law subject_ko=민법 seq=16 number=16
%% source: Q-Net 감정평가사 기출문제 (cst00309) artlSeq=5052578 file=raw/1008/5052578/2012년 제23회 감정평가사 제1차시험 기출문제(1교시 A형).pdf
%% conversion: kordoc 4.24.51 -> tools/gen1.py (auto); review=proofread-sonnet; flags=none
\begin{question}{16}
통정허위표시에 있어 선의의 제3자가 될 수 없는 자는? (다툼이 있으면 판례에 의함)
\begin{choices}
  \item 가장채권 양수인의 포괄승계인
  \item 가장행위에 기한 저당권부 채권을 가압류한 자
  \item 가장양수인으로부터 매매계약에 의한 소유권이전등기청구권 보전을 위한 가등기를 경료받은 자
  \item 가장저당권의 실행으로 경락받은 자
  \item 가장양수인으로부터 목적부동산을 전득한 자
\end{choices}
\end{question}
\answerall
\answernote{정답표에 1,2,3,4,5 모두 정답으로 표기(전원 정답)}
